\section{Results}
\label{sec:results}

\subsection{Data harmonisation and SPARK-only target}

The two sumstats files are on different genome builds (iPSYCH--PGC on
GRCh37, the Stein-Lab meta-analysis on GRCh38). A na\"{i}ve
\texttt{chr:pos} inner-join therefore matched only $\sim 73$\,k of the
$\sim 9$\,M SNPs in either file; the rsID-based join described in
Section~\ref{sec:methods} recovered $7\,413\,400$ shared SNPs, of which
$6\,309\,737$ survived QC (37\,101 in MHC). The harmonised file has
median $|Z_{\mathrm{iPSYCH}}| = 0.720$ and median
$|Z_{\mathrm{meta}}| = 0.724$, both consistent with mild polygenic
inflation.

Round-trip validation of Eq.~\ref{eq:sparkz} reproduces the original
$Z_{\mathrm{meta}}$ to floating-point precision (max error
$8.9 \times 10^{-16}$), confirming that the inversion is exact once
$N_{\mathrm{iPSYCH}}$ and $N_{\mathrm{meta}}$ are specified
correctly. The derived SPARK-only Z-scores have median absolute value
$0.715$ and a Pearson correlation of $0.035$ with the discovery Z, as
expected for two effectively independent samples that agree on the
underlying true effects but differ in their large per-SNP measurement
noise.

The Manhattan and QQ plots of the discovery sumstats
(Fig.~\ref{fig:manhattan},~\ref{fig:qq}) reproduce the genome-wide
significant peaks reported by Grove et al.~\cite{grove2019} and exhibit
the characteristic upward-bending QQ tail consistent with polygenic
inflation rather than confounding.

\begin{figure}[t]
  \centering
  \IfFileExists{figures/fig_manhattan.pdf}{%
    \includegraphics[width=\linewidth]{fig_manhattan.pdf}%
  }{\fbox{\parbox{0.95\linewidth}{\centering Run \texttt{make all} in
       \texttt{code/} to generate \texttt{fig\_manhattan.pdf}.}}}
  \caption{Manhattan plot of the iPSYCH--PGC discovery sumstats
    (6\,309\,737 SNPs after harmonisation and QC).
    The horizontal red line marks $p = 5 \times 10^{-8}$.}
  \label{fig:manhattan}
\end{figure}

\begin{figure}[t]
  \centering
  \IfFileExists{figures/fig_qq.pdf}{%
    \includegraphics[width=0.85\linewidth]{fig_qq.pdf}%
  }{\fbox{\parbox{0.85\linewidth}{\centering Run \texttt{make all} to
       generate \texttt{fig\_qq.pdf}.}}}
  \caption{QQ plot of $-\log_{10}(p)$ for the iPSYCH--PGC discovery
    sumstats. The diagonal reference line marks the null expectation.}
  \label{fig:qq}
\end{figure}

\subsection{Held-out predictive accuracy}

We evaluate every method on chr~21--22 (170\,039 test SNPs after
removing MHC). Table~\ref{tab:eval} reports Pearson $r$ between the
method-predicted Z-score and the SPARK-only Z, with bootstrapped 95\,\%
confidence intervals (1000 SNP-level resamples), the AUC for the binary
task $|Z_{\mathrm{SPARK}}| > 1.96$, and the calibration slope.
Fig.~\ref{fig:bars} renders the same numbers visually.

\begin{table}[t]
  \centering
  \scriptsize
  \setlength{\tabcolsep}{3pt}
  \begin{tabular}{lrrr}
    \toprule
    Method & $n$ & $r$ [95\,\% CI] & AUC \\
    \midrule
    \textbf{ml\_xgb}     & 170\,039 & \textbf{0.058} [0.053, 0.063] & 0.488 \\
    ml\_mlp              & 170\,039 & 0.048 [0.043, 0.053] & 0.499 \\
    PRS-CS-EB            & 170\,039 & 0.041 [0.036, 0.046] & 0.499 \\
    GNN (no graph)       & 170\,039 & 0.037 [0.033, 0.042] & 0.499 \\
    \midrule
    P+T $p \!\leq\! 1.0$     &      213 & $-0.046$ [$-0.18$, $0.09$] & 0.587 \\
    P+T $p \!\leq\! 5e\!-\!2$ & 191 & $-0.045$ [$-0.17$, $0.09$] & 0.537 \\
    P+T $p \!\leq\! 10^{-3}$ &       37 & $-0.097$ [$-0.36$, $0.20$] & 0.578 \\
    P+T $p \!\leq\! 10^{-5}$ &        2 & ---                          & --- \\
    P+T $p \!\leq\! 5e\!-\!8$ &  0 & ---                          & --- \\
    \bottomrule
  \end{tabular}
  \caption{Held-out test-fold metrics (chr~21--22, MHC excluded). Rows
    above the inner rule are evaluated on every test SNP; P+T rows
    include only the SNPs that survived clumping at the listed
    threshold (\texttt{---} where $n \!<\! 3$ makes the metric
    meaningless).}
  \label{tab:eval}
\end{table}

\begin{figure}[t]
  \centering
  \IfFileExists{figures/fig_method_bars.pdf}{%
    \includegraphics[width=\linewidth]{fig_method_bars.pdf}%
  }{\fbox{\parbox{0.95\linewidth}{\centering Run \texttt{make all} to
       generate \texttt{fig\_method\_bars.pdf}.}}}
  \caption{Held-out predictive accuracy. Left: Pearson $r$ between the
    method-predicted Z and the SPARK-only Z on chr~21--22, with
    bootstrapped 95\,\% CIs. Right: AUC for the binary task
    $|Z_{\mathrm{SPARK}}| > 1.96$.}
  \label{fig:bars}
\end{figure}

\begin{figure}[t]
  \centering
  \IfFileExists{figures/fig_method_scatter.pdf}{%
    \includegraphics[width=\linewidth]{fig_method_scatter.pdf}%
  }{\fbox{\parbox{0.95\linewidth}{\centering Run \texttt{make all} to
       generate \texttt{fig\_method\_scatter.pdf}.}}}
  \caption{Per-method predicted (x-axis) versus observed SPARK-only
    Z-score (y-axis) on chr~21--22, hex-binned. Each in-text $r$ value
    is annotated.}
  \label{fig:scatter}
\end{figure}

The four whole-genome methods produce a clear ordering. The
\textbf{XGBoost re-weighter} attains the highest Pearson $r = 0.058$
(95\,\% CI $[0.053, 0.063]$), beating the empirical-Bayes baseline by
0.017 and the discovery-only \textit{GNN-no-graph} fallback by 0.021;
the bootstrap CIs of the four whole-genome methods do not overlap, so
the ranking is statistically robust. The \textbf{MLP re-weighter} comes
second ($r = 0.048$). Both ML methods improve over the empirical-Bayes
baseline despite using only the discovery Z, $|Z|$, INFO score, and
LD-window summaries -- the gene-level functional annotations sketched
in Section~\ref{sec:methods} were not available in this run because the
required reference files (Gencode, GTEx, ENCODE, SFARI Gene, STRING)
are gated by separate download steps; we expect the gain to widen
materially once those features are added.

The \textbf{P+T baselines} are dominated on this evaluation. Even the
most permissive threshold ($p \leq 1$) retains only 213 chr~21--22
SNPs after clumping, and at $p \leq 5\!\times\!10^{-8}$ no chr~21--22
SNP survives -- a direct illustration of P+T's known limitation for a
highly polygenic trait. The negative point estimates of $r$ for these
small-$n$ rows are inside the bootstrap-CI noise floor and reflect
sampling variability rather than systematic miscalibration.

The AUC for the binary "significance discrimination'' task hovers
near 0.5 for the whole-genome methods, which is expected: very few
chr~21--22 SNPs reach $|Z_{\mathrm{SPARK}}| > 1.96$ at
$N_{\mathrm{SPARK}} = 12{,}444$, so this metric has little power on
this hold-out and should not be over-interpreted.

\subsection{Calibration}

Calibration slopes (regressing observed-$Z$ on predicted-$Z$ on the test
fold) tell a complementary story.\ The MLP re-weighter is the closest
to the ideal slope of $1.0$ (slope $0.90$), the XGBoost re-weighter
slightly over-shrinks ($1.83$), and the empirical-Bayes baseline is
under-confident ($0.058$, predicted $Z$ values are an order of
magnitude smaller than observed). This pattern -- ML models that
directly target the held-out Z produce better-calibrated posteriors
than fixed-form Bayesian shrinkage -- is a useful additional argument
for the proposed reweighter approach.

\subsection{Stratified analyses}

\texttt{eval\_strata.tsv} reports Pearson $r$ within imputation-INFO
bins. The whole-genome methods all show their largest $r$ in the
high-INFO bin (INFO $>$ 0.95: $r \approx 0.06$--$0.08$), and the
relative ordering of methods is preserved. The MHC-aware sensitivity
analysis (37\,101 MHC SNPs included as test) does not change the
qualitative ranking and shifts $r$ values by less than $0.005$
(detailed in the supplementary parquet output).

\subsection{Biological enrichment}

The enrichment analysis described in Section~\ref{sec:methods} is gated
by the optional reference annotation files. Because the Gencode v44
gene BED, SFARI Gene catalogue, STRING v12 edges, and GTEx eQTL tracks
were not available locally for this run, the SNP-to-gene mapping was
empty and the per-method gene-set Fisher tests in
\texttt{enrichment.tsv} therefore record a \texttt{missing-ref} status
rather than $p$-values. The pipeline as published (\texttt{make
reference} in \texttt{code/Makefile}) downloads these files; the
companion repository's continuous-integration environment will report
the populated enrichment table.
